% 若编译失败，且生成 .synctex(busy) 辅助文件，可能有两个原因：
% 1. 需要插入的图片不存在：Ctrl + F 搜索 'figure' 将这些代码注释/删除掉即可
% 2. 路径/文件名含中文或空格：更改路径/文件名即可

% --------------------- 文章宏包及相关设置 --------------------- %
% >> ------------------ 文章宏包及相关设置 ------------------ << %
% 设定文章类型与编码格式
\documentclass[UTF8]{article}		
\input{../.config/config_for_LinearCircuitExperiment.tex}



%%%%%%%%%%%%%%%%%%%%%%%%%%%%%%%%%%%%%%%%%%%%%%%%%%%%%%%%%%%%%%%%
% 仅需修改页眉、实验名称、实验日期
%%%%%%%%%%%%%%%%%%%%%%%%%%%%%%%%%%%%%%%%%%%%%%%%%%%%%%%%%%%%%%%%


%%%%%%%%%%%%%%%%%% 1. 修改页眉内容 %%%%%%%%%%%%%%%%%%
\rhead{LCE-09 比较器 (2025.06.06, 丁毅)}

% 开始编辑文章
\begin{document}
\begin{center}\large
    \vspace*{-0.8cm}
    \noindent{\huge\bfseries《\ \ 线\ \ 性\ \ 电\ \ 路\ \ 实\ \ 验\ \ \ 》\ \ 实\ \ 验\ \ 报\ \ 告 }
    \\\vspace{0.1cm}
    \noindent{
    {\bfseries 
%
%%%%%%%%%%%%%%%%%% 2. 修改实验名称 %%%%%%%%%%%%%%%%%%
    实验名称：\uline{\hspace{2.2cm} 比较器 \hspace{2.2cm}}
%
    }\hspace{0.4cm}
    指导教师：\uline{\hspace{0.8cm}王东雷\ \ \ \  \ df4dac@sina.com     \hspace{0.8cm}}
    }
    \\\vspace{0.1cm}
    \noindent
    {
    姓名：\uline{\,\,\,丁毅\,\,\,}\hspace{0.2cm}
    学号：\uline{\,\,\,{ 2023K8009908031}\,\,\,}\hspace{0.2cm}
    班级/专业：\uline{\,\,\,{2308/电子信息}\,\,\,}\hspace{0.2cm}
    分组序号：\uline{\,\,\,{2-06}\,\,\,}
    }
    \\\vspace{0.1cm}
    \noindent{
%
%%%%%%%%%%%%%%%%%% 3. 修改实验日期 %%%%%%%%%%%%%%%%%%
    实验日期：\uline{\,\,{2025.06.06}\,\,}\hspace{0.2cm}
%
    实验地点：\uline{\,\,\,教学楼{ 607}\,\,\,}\hspace{0.2cm}
    是否调课/补课：\uline{\hspace{0.26cm}否 \hspace{0.26cm}}\hspace{0.2cm}
    成绩：\uline{\hspace{2cm}}
    }
\end{center}
\vspace{-0.4cm}
\noindent\rule{\textwidth}{0.075em}   % 分割线
\vspace{-1.0cm}

% 生成目录
\setcounter{tocdepth}{3}  % 目录深度为 2（不显示 subsubsection）
\noindent\tableofcontents\thispagestyle{fancy}   % 显示页码、页眉等

% ------------------------ 文章信息区 ------------------------ %
% ------------------------ 文章信息区 ------------------------ %



%%%%%%%%%%%%%%%%%%%%%%%%%%%%%%%%%%%%%%%%%%%%%%%%%%%%%%%%%%%%%%%%%%%%%%%%%%%%%%%%%
%%%%%%%%%%%%%%%%%%%%%%%%%%%%%%%%% 下面是正文内容 %%%%%%%%%%%%%%%%%%%%%%%%%%%%%%%%%
%%%%%%%%%%%%%%%%%%%%%%%%%%%%%%%%%%%%%%%%%%%%%%%%%%%%%%%%%%%%%%%%%%%%%%%%%%%%%%%%%

\section{实验目的}
\begin{enumerate}
\item 理解比较器的主要参数、比较器与运放的差异；
\item 测量比较器时间延迟等参数，观察滞回特性对电路的影响；
\item 设计电压监测比较器电路；
\item 了解常见比较器 LM293P 和 LM311P。
\end{enumerate}


\section{实验仪器}

\begin{enumerate}
    \item 数字万用表： Unit UT61E (C190241394)
    \item 数字示波器： RIGOL 200MSO2202A (DS2F192200361)
    \item 信号发生器： GWINSTEK AFG-22225 (GER910370)
    \item 数字直流电源： GWINSTEK GPD-3303S (GES813705) 
    \item 多功能数字测量仪： 
    % 这里是链接
    \href{https://digilent.com/reference/test-and-measurement/analog-discovery/start
    }{ % 这里是文字
    Analog Discovery 1
    } 
    (D704387)
    \item 其它：比较器 LM311P 、LM293P 等，电阻、电容、二极管等元件，以及面包板，导线等
\end{enumerate}

\section{实验原理}

\subsection{比较器与运放的差异}

比较器工作于开环或正反馈（滞回）模式，且内部无需频率补偿，运放一般工作在闭环负反馈模式。相近的电路，比较器比运放上升速度快。运放可当作低速比较器，比较器用作运放会振荡（无频率补偿）。比较器工作时会有较高的差模电压。一般后接数字电路或控制电路，输出端一般还有地线端子，运放没有地线端子。双极型比较器常用 OC (Open-Collector) 输出结构，结合上拉电阻实现逻辑电平的输出，缺点是功耗会较大，且驱动能力不足。现在的新型 CMOS 比较器基本都实现了 rail-to-rail input/output ，可以与逻辑电平兼容，且输出采用推挽结构，省掉了上拉电阻的同时，提升了驱动能力，降低了功耗。

\subsection{由滞回比较器构成带有滞回的窗口比较器}

采用两个输入方式不同的电压比较器组合可以形成窗口比较器，利用窗口比较器进行某些物理量的监测的逻辑类似于与门。加入滞回的目的是：由于比较器是高速器件，放在电路中进行电平转换时会在转换点附近产生高频振荡，加入滞回以后可以利用滞回特性消除这一振荡，从而使得监测更加准确。

\section{实验内容及实验结果}

\vspace*{-3mm}
\begin{redbox}
    若无特别说明，后续实验均默认示波器为 dc coupling 模式。
\end{redbox}

\subsection{搭建并测量无滞回比较器}


\begin{figure}[H]\centering
    \includegraphics[width=0.6\columnwidth]{LCE-09-比较器/assets/滞回比较器.pdf}
    \caption{滞回比较器原理图 (对于无滞回比较器，只需去掉 RF2, 并将 RF1 替换为导线，使得 Pin 2 直接接地)}
\end{figure}


先不接入滞回，用示波器观察输出波形如下：

\begin{figure}[H]\centering
    \includegraphics[width=0.95\columnwidth]{LCE-09-比较器/assets/无滞回 1.png}
    \caption{LM311P 无滞回比较器的输入输出波形; CH1 (orange) 为输入信号，CH2 (blue) 为输出信号}
\end{figure}
\begin{figure}[H]\centering
    \includegraphics[width=0.95\columnwidth]{LCE-09-比较器/assets/无滞回 2.png}
    \caption{LM311P 无滞回比较器的上升沿波形; CH1 (orange) 为输入信号，CH2 (blue) 为输出信号} 
\end{figure}
\begin{figure}[H]\centering
    \includegraphics[width=0.95\columnwidth]{LCE-09-比较器/assets/无滞回 3.png}
    \caption{LM311P 无滞回比较器的下降沿波形; CH1 (orange) 为输入信号，CH2 (blue) 为输出信号}
\end{figure}

不妨观察一下 RF1 未短路时的输出波形 (Pin 2 通过 RF1 连接到 GND)，翻转时出现明显振荡：

\begin{figure}[H]\centering
    \includegraphics[width=\columnwidth]{LCE-09-比较器/assets/RF1 1.png}
    \caption{LM311P 无滞回比较器的上升沿波形 (RF1 未短路); CH1 (orange) 为输入信号，CH2 (blue) 为输出信号}
    \label{fig: RF1 1}
\end{figure}

\begin{figure}[H]\centering
    \includegraphics[width=\columnwidth]{LCE-09-比较器/assets/RF1 2.png}
    \caption{LM311P 无滞回比较器的下降沿波形 (RF1 未短路); CH1 (orange) 为输入信号，CH2 (blue) 为输出信号}
    \label{fig: RF1 2}
\end{figure}

\newpage
\subsection{搭建并测量有滞回比较器}

正常接入 RF1 和 RF2 后，输入输出波形如下：

\begin{figure}[H]\centering
    \includegraphics[width=0.87\columnwidth]{LCE-09-比较器/assets/有滞回 1.png}
    \caption{LM311P 有滞回比较器的输入输出波形; CH1 (orange) 为输入信号，CH2 (blue) 为输出信号}
\end{figure}
\begin{figure}[H]\centering
    \includegraphics[width=0.87\columnwidth]{LCE-09-比较器/assets/有滞回 2.png}
    \caption{LM311P 有滞回比较器的上升沿波形; CH1 (orange) 为输入信号，CH2 (blue) 为输出信号}
\end{figure}
\begin{figure}[H]\centering
    \includegraphics[width=0.87\columnwidth]{LCE-09-比较器/assets/有滞回 3.png}
    \caption{LM311P 有滞回比较器的下降沿波形; CH1 (orange) 为输入信号，CH2 (blue) 为输出信号}
\end{figure}
\begin{figure}[H]\centering
    \includegraphics[width=0.87\columnwidth]{LCE-09-比较器/assets/有滞回 4.png}
    \caption{下降沿波形细节; CH1 (orange) 为输入信号，CH2 (blue) 为输出信号}
\end{figure}


\subsection{搭建并测量窗口比较器}

要求供电电压的正常范围是 10 V $\sim$ 15 V, 超出范围警报灯亮起。过压比较器采用同相输入的
方式、欠压比较器采用反相输入的方式。输出接上拉电阻后接一个二极管，后续电路接入一个发
光二极管和一个限流电阻，在限制范围之外时会有一个比较器输出高电平使二极管导通，从而使
发光二极管获得足够电压后发光，示意电压异常。参考端利用 
\href{https://digilent.com/reference/test-and-measurement/analog-discovery/start
}{ % 这里是文字
    Analog Discovery 1
} 
输出 5 V 作为参考电压 $V_{REF}$ 。


电路原理图及实物图如下所示：

\begin{figure}[H]\centering
\begin{subfigure}[b]{0.5\columnwidth}\centering
    \includegraphics[height=165pt]{LCE-09-比较器/assets/窗口比较器.pdf}
    \caption{原理图}
\end{subfigure}\hfill
\begin{subfigure}[b]{0.5\columnwidth}\centering
    \includegraphics[height=165pt]{LCE-09-比较器/assets/image.png}
    \caption{实物图}
\end{subfigure}
\caption{用两个滞回比较器搭建具有滞回功能的窗口比较器}
\end{figure}

调整电源电压 VCC, 得到阈值电压数据如下：

\begin{table}[H]\centering
    %\renewcommand{\arraystretch}{1.5} % 调整行间距
    %\setlength{\tabcolsep}{1.5mm} % 调整列间距
    \caption{窗口比较器的阈值电压数据 (0 表示灯灭, 1 表示灯亮)}
\begin{tabular}{cccccccccc}\toprule
    $V_{H, 0 \to 1}$ & $V_{H, 1 \to 0}$ &  $V_{L, 0 \to 1}$ & $V_{L, 1 \to 0}$ \\
    \midrule
    15.003 V & 14.939 V & 9.954 V & 9.978 V  \\
    \bottomrule
\end{tabular}
\end{table}

检测电压的阈值偏差和滞回限度都比较小，可以认为达到了设计要求。

\section{思考题}

\subsection{比较器使用上有哪些注意事项？}

比较器内部没有频率补偿，在电平翻转时很容易产生振荡。例如我们在 Figure \ref{fig: RF1 1} 和 Figure \ref{fig: RF1 2} 中可以看到，电平翻转过程中，上升沿和下降沿均出现了明显振荡。

\subsection{运放做比较器使用，有什么优缺点？}

运放通常有足够的相位裕度，因此在电平翻转时输出通常更稳定，不容易产生振荡。缺点是普通运放的响应时间 (response time, 或建立时间 setting time) 较长，输出摆率低，无法处理高速时钟信号。

\subsection{实验中是否遇到异常现象，如何解决？}

在搭建窗口比较器时，我们遇到了“注入正弦信号，但输出电平没有翻转”的情况，这是输入信号幅度不够导致的。输入信号幅度并不足以让比较器通过滞回点，因此输出电平没有翻转，增大输入信号幅度后，问题解决。





\section{开关电源设计预告}


\begin{figure}[H]\centering
\begin{subfigure}[b]{0.5\columnwidth}\centering
    \includegraphics[height=210pt]{LCE-09-比较器/assets/开关 1.png}
\end{subfigure}\hfill
\begin{subfigure}[b]{0.5\columnwidth}\centering
    \includegraphics[height=210pt]{LCE-09-比较器/assets/开关 2.png}
\end{subfigure}
\caption{开关电源设计预告}
\end{figure}


















% 附录
\section*{附录 A\hspace*{20pt} 原始数据记录表}
\addcontentsline{toc}{section}{附录 A\hspace*{6pt} 原始数据记录表} 
\thispagestyle{fancy} 
\begin{figure}[H]\centering
    
\includegraphics[width=\columnwidth]{LCE-09-比较器/assets/data.png}
\end{figure}

% 附录
\newpage
\vspace*{\fill}\begin{center}\Huge{\bfseries 
    附录 B\hspace*{20pt} 实验预习报告
}\end{center}\addcontentsline{toc}{section}{附录 B\hspace*{6pt} 实验预习报告} 
\vspace*{\fill}
\thispagestyle{fancy} 
\includepdf[pages={-}]{LCE-09-比较器/preview/LCE-09 (preview report).pdf}


% 附录
%\section*{附录 C \hspace*{20pt} Matlab Codes}
%\addcontentsline{toc}{section}{附录 C \hspace*{6pt} Matlab Codes} 
%\thispagestyle{fancy} 
%\lstinputlisting{d:/a_RemoteRepo/GH.MatlabCodes/本科课程代码/基础物理实验/Ex_9/Ex_9_analysis_mfile.m}



\end{document}

% VScode 常用快捷键：

% F2:                       变量重命名
% Ctrl + Enter:             行中换行
% Alt + up/down:            上下移行
% 鼠标中键 + 移动:           快速多光标
% Shift + Alt + up/down:    上下复制
% Ctrl + left/right:        左右跳单词
% Ctrl + Backspace/Delete:  左右删单词    
% Shift + Delete:           删除此行
% Ctrl + J:                 打开 VScode 下栏(输出栏)
% Ctrl + B:                 打开 VScode 左栏(目录栏)
% Ctrl + `:                 打开 VScode 终端栏
% Ctrl + 0:                 定位文件
% Ctrl + Tab:               切换已打开的文件(切标签)
% Ctrl + Shift + P:         打开全局命令(设置)

% Latex 常用快捷键：

% Ctrl + Alt + J:           由代码定位到PDF


